% Options for packages loaded elsewhere
\PassOptionsToPackage{unicode}{hyperref}
\PassOptionsToPackage{hyphens}{url}
\PassOptionsToPackage{dvipsnames,svgnames,x11names}{xcolor}
\documentclass[
  10pt,
  fontset=fandol]{ctexart}
\usepackage{xcolor}
\usepackage[margin=16mm]{geometry}
\usepackage{amsmath,amssymb}
\setcounter{secnumdepth}{-\maxdimen} % remove section numbering
\usepackage{iftex}
\ifPDFTeX
  \usepackage[T1]{fontenc}
  \usepackage[utf8]{inputenc}
  \usepackage{textcomp} % provide euro and other symbols
\else % if luatex or xetex
  \usepackage{unicode-math} % this also loads fontspec
  \defaultfontfeatures{Scale=MatchLowercase}
  \defaultfontfeatures[\rmfamily]{Ligatures=TeX,Scale=1}
\fi
\usepackage{lmodern}
\ifPDFTeX\else
  % xetex/luatex font selection
\fi
% Use upquote if available, for straight quotes in verbatim environments
\IfFileExists{upquote.sty}{\usepackage{upquote}}{}
\IfFileExists{microtype.sty}{% use microtype if available
  \usepackage[]{microtype}
  \UseMicrotypeSet[protrusion]{basicmath} % disable protrusion for tt fonts
}{}
\makeatletter
\@ifundefined{KOMAClassName}{% if non-KOMA class
  \IfFileExists{parskip.sty}{%
    \usepackage{parskip}
  }{% else
    \setlength{\parindent}{0pt}
    \setlength{\parskip}{6pt plus 2pt minus 1pt}}
}{% if KOMA class
  \KOMAoptions{parskip=half}}
\makeatother
\setlength{\emergencystretch}{3em} % prevent overfull lines
\providecommand{\tightlist}{%
  \setlength{\itemsep}{0pt}\setlength{\parskip}{0pt}}
\usepackage{amsmath,amssymb,mathtools,needspace}
\xeCJKDeclareCharClass{CJK}{"2032,"0400->"04FF,"2160->"216B,"2460->"2473,"25A0,"25C7,"25CB,"25CF}
\IfFontExistsTF{Arial Unicode MS}{\xeCJKsetup{AutoFallBack=true}\setCJKfallbackfamilyfont{\CJKrmdefault}{Arial Unicode MS}}{}
\setcounter{secnumdepth}{0}
\setlength{\emergencystretch}{2em}
\allowdisplaybreaks
\usepackage{bookmark}
\IfFileExists{xurl.sty}{\usepackage{xurl}}{} % add URL line breaks if available
\urlstyle{same}
\hypersetup{
  pdftitle={其他常用的正交多项式},
  colorlinks=true,
  linkcolor={Maroon},
  filecolor={Maroon},
  citecolor={Blue},
  urlcolor={Blue},
  pdfcreator={LaTeX via pandoc}}

\title{其他常用的正交多项式}
\author{}
\date{}

\begin{document}
\maketitle

\subsubsection{3.4.4 其他常用的正交多项式}\label{ux5176ux4ed6ux5e38ux7528ux7684ux6b63ux4ea4ux591aux9879ux5f0f}

一般地说, 如果区间 \([a,b]\) 及权函数 \(\rho(x)\) 不同, 则得到的正交多项式也不同. 除上述两种最重要的正交多项式外, 下面再给出三种较常用的正交多项式.

\paragraph{1. 第二类 Chebyshev 多项式}\label{ux7b2cux4e8cux7c7b-chebyshev-ux591aux9879ux5f0f}

在区间 \([-1,1]\) 上带权 \(\rho(x) = \sqrt{1-x^2}\) 的正交多项式称为第二类 Chebyshev 多项式, 其表达式为

\[
U_n(x) = \frac{\sin[(n+1)\arccos x]}{\sqrt{1-x^2}}. \tag{3.4.10}
\]

由 \(x=\cos\theta\) 可得

\[
\int_{-1}^{1} U_n(x) U_m(x) \sqrt{1-x^2} \mathrm{d}x = \int_{0}^{\pi} \sin(n+1)\theta \sin(m+1)\theta \mathrm{d}\theta
\]

\[
= \left\{ \begin{array}{l l} 0, & m \neq n, \\ \frac{\pi}{2}, & m=n, \end{array} \right.
\]

\begin{quote}
校注（agent 补充，记号约定）：式 (3.4.9) 后原文规定 \(T_0=1/2\)，而表 3.2 首行原印 \(1=T_0\)；两处均忠实保留。式 (3.4.9) 的求和使用临时的半权常数项约定，表 3.2 则与前文通常定义 \(T_0=1\) 一致。使用时须区分这两种约定。
\end{quote}

\href{../../source-images/pdf-076.jpeg}{查看原页}

即 \(\{U_n(x)\}\) 是区间 \([-1,1]\) 上带权 \(\rho(x)=\sqrt{1-x^2}\) 的正交多项式族。还可得到递推关系式

\[
U_0(x)=1,\quad U_1(x)=2x,
\]

\[
U_{n+1}(x)=2xU_n(x)-U_{n-1}(x)\quad(n=1,2,\cdots).
\]

\paragraph{2. Laguerre 多项式}\label{laguerre-ux591aux9879ux5f0f}

在区间 \([0,\infty)\) 上带权 \(\rho(x)=\mathrm e^{-x}\) 的正交多项式称为 Laguerre 多项式，其表达式为

\[
L_n(x)=\mathrm e^n\frac{\mathrm d^n}{\mathrm dx^n}(x^n\mathrm e^{-x}).\tag{3.4.11}
\]

它也具有正交性质

\[
\int_0^\infty \mathrm e^{-x}L_n(x)L_m(x)\,\mathrm dx=
\begin{cases}
0,&m\ne n,\\
(n!)^2,&m=n
\end{cases}
\]

和递推关系

\[
L_0(x)=1,\quad L_1(x)=1-x,
\]

\[
L_{n+1}(x)=(1+2n-x)L_n(x)-n^2L_{n-1}(x)\quad(n=1,2,\cdots).
\]

\paragraph{3. Hermite 多项式}\label{hermite-ux591aux9879ux5f0f}

在区间 \((-\infty,\infty)\) 上带权 \(\rho(x)=\mathrm e^{-x^2}\) 的正交多项式称为 Hermite 多项式，其表达式为

\[
H_n(x)=(-1)^n\mathrm e^{x^2}\frac{\mathrm d^n}{\mathrm dx^n}(\mathrm e^{-x^2}).\tag{3.4.12}
\]

它满足正交关系

\[
\int_{-\infty}^{+\infty}\mathrm e^{-x^2}H_m(x)H_n(x)\,\mathrm dx=
\begin{cases}
0,&m\ne n,\\
2^nn!\sqrt\pi,&m=n,
\end{cases}
\]

并有递推关系

\[
H_0(x)=1,\quad H_1(x)=2x,
\]

\[
H_{n+1}(x)=2xH_n(x)-2nH_{n-1}(x)\quad(n=1,2,\cdots).
\]

\begin{center}\rule{0.5\linewidth}{0.5pt}\end{center}

\subsection{相关原页校注（agent 补充）}\label{ux76f8ux5173ux539fux9875ux6821ux6ce8agent-ux8865ux5145}

以下校注来自本节覆盖的原页，可能也涉及同页相邻小节；原文照录与校正说明分开保留。

\subsubsection{原书 PDF 页 76 的校注}\label{ux539fux4e66-pdf-ux9875-76-ux7684ux6821ux6ce8}

\href{../pages/pdf-076.md}{完整原页转录} · \href{../../source-images/pdf-076.jpeg}{原页图像}

校注记录：ch03b-E001。

\begin{quote}
校注（agent 补充）：式 (3.4.11) 的原扫描明确印为 \(\mathrm e^n\)，已忠实保留；这是疑似原书错误，不是 OCR 误读。取 \(n=0\)，原式给出 \(L_0(x)=\mathrm e^{-x}\)，与紧接着印出的 \(L_0(x)=1\) 矛盾。相容的 Rodrigues 表达式前因子应为 \(\mathrm e^x\)。见 \href{../assets/ch03b/p076-laguerre-formula.png}{原式放大裁切}。
\end{quote}

\subsection{本节覆盖原页的校注记录}\label{ux672cux8282ux8986ux76d6ux539fux9875ux7684ux6821ux6ce8ux8bb0ux5f55}

下列记录按原页关联，含同页相邻小节的校注；计算时同时读取上文校注和完整原页。

\begin{itemize}
\tightlist
\item
  \texttt{ch03b-E001}：原书 PDF 页 76
\end{itemize}

\end{document}
